\paragraph{Relation to the source constructions.}
The period matrix, two affine finite fillings, logarithmic collar maps
and untranslated cusp fan are the original construction
\cite[Sections~3--6]{AC2026}.  The same local pieces in the two
versions of \cite{DuEtAl2026v1,DuEtAl2026v2} are identified by the
parameter translation $a=c+\delta$ of \cref{night:P5.7}; in particular
$\Omega_a=\Pi_c$, with the ordered lattice, vectors and fan retained.
The scalar cocycle below is precisely $f=-\eta_{\rm DGKX}$ on every
displayed generating arrow in
\cite[Section~6.4, the scalar cocycle preceding Lemma~6.14]{DuEtAl2026v1}.
The minus sign is the stated line-coordinate convention.  The
exponential quotient comparison with Engel's elliptic-surface
presentation is the map already given in
\cref{fam:src:B.1.1}.

Conversion of a holomorphic cocycle to a Dolbeault form by smooth
weights is the usual \v Cech--Dolbeault construction
\cite[Theorem~3.1 and proof, p.~10]{AbateBracciSuwaTovena2010Atiyah}.
On two charts its formula applied to $(0,0,c_{01}/12)$ is
$-\bar\partial w_0\wedge c_{01}/12
=\bar\partial(w_1c_{01})/12$, exactly the convention used here;
replace $c_{01}/12$ by $f_{01}$ for $G$.
The proof supplies the actual local theta and Cartier lines, their
annular holomorphic connections, the resulting one-form cocycle,
and a specified locally finite coordinate weight system with every
derivative evaluated.  Thus the general cohomological conversion and
the scalar input are credited together with the original family data.
